% Lösungen Einheit 2 von 3 – Die Nullhypothese wählen: aus der Sicht des Entscheiders (zusammenbau v0.1)
\einheitenkopf{Einheit 2 von 3 · Die Nullhypothese wählen: aus der Sicht des Entscheiders}

\erg{11}{a) Fehler erster Art \quad b) Auf das teurere Gummi wird nur umgestellt, wenn die Nullhypothese abgelehnt wird, also nur bei nachgewiesener Verbesserung; die Wahrscheinlichkeit, irrtümlich teuer umzustellen, obwohl sich nichts verbessert hat, ist als Fehler erster Art durch das Signifikanzniveau begrenzt. \quad c) Abgeschaltet wird nur, wenn die Nullhypothese abgelehnt wird, also nur bei nachgewiesenem Rückgang; das Risiko, den Algorithmus irrtümlich abzuschalten, obwohl die Zufriedenheit nicht gesunken ist, ist durch das Signifikanzniveau begrenzt. \quad d) Der Händler will vermeiden, eine schlechte Lieferung irrtümlich anzunehmen: Er nimmt sie nur an, wenn die Nullhypothese abgelehnt wird, und die Wahrscheinlichkeit, eine Lieferung mit zu vielen defekten Geräten irrtümlich anzunehmen, ist dann als Fehler erster Art durch das Signifikanzniveau begrenzt.}
\erg{12}{a) Begrenzt ist nur der Fehler erster Art: Eine schlechte Lieferung anzunehmen heißt hier, die zutreffende Nullhypothese abzulehnen – genau deshalb ist dieses Risiko durch das Signifikanzniveau begrenzt. \quad b) Die Entscheidungsregel wird so festgelegt, dass eine Ablehnung bei zutreffender Nullhypothese höchstens mit der Wahrscheinlichkeit des Niveaus eintritt; wie oft eine falsche Nullhypothese beibehalten wird, hängt vom unbekannten wahren Anteil ab und wird nicht festgelegt. \quad c) Nullhypothese: Höchstens 25\,\% der Anwohner fühlen sich gestört. Die Wand wird nur bei Ablehnung gebaut; das Risiko, sie irrtümlich zu bauen, ist durch das Signifikanzniveau begrenzt.}
